Open de webbrowser die op de computer beschikbaar is.

Zoek informatie over Linux. Gebruik bijvoorbeeld de website van Debian of de documentatie van KDE.

Voer de volgende opdrachten uit:
\begin{enumerate}
\item Zoek op internet naar informatie over Debian.
\item Open een nieuw tabblad.
\item Zoek een webpagina met uitleg over KDE Plasma.
\item Open een nieuw tabblad.
\item Zoek een webpagina met uitleg over LibreOffice.
\end{enumerate}

Beantwoord de vragen:
\begin{enumerate}
\item Welke browser heb je gebruikt?
\item Wat is het verschil tussen een zoekmachine en een webbrowser?
\item Wat is een bladwijzer en waarvoor gebruik je die?
\item Noteer de titel en het webadres van twee gebruikte informatiebronnen.
\item Waarom is het belangrijk om informatie uit meerdere bronnen te controleren?
\end{enumerate}

Let op: download geen programma's en voer geen bestanden uit die je op Internet vindt.

Sla de antwoorden op in een bestand met de naam {Linux\_Informatie}.

